\chapter{Parry}

\lettrine{W}{hilst} the unknown was viewing these lights with  interest, and lending an ear to the various noises, Master Cropole entered his  apartment, followed by two attendants, who laid the cloth for his meal.

\zz
The stranger did not pay them the least attention; but Cropole approaching him  respectfully, whispered, <Monsieur, the diamond has been valued.>

<Ah!> said the traveller. <Well?>

<Well, monsieur, the jeweller of S. A. R. gives two hundred and eighty  pistoles for it.>

<Have you them?>

<I thought it best to take them, monsieur; nevertheless, I made it a  condition of the bargain, that if monsieur wished to keep his diamond, it  should be held till monsieur was again in funds.>

<Oh, no, not at all: I told you to sell it.>

<Then I have obeyed, or nearly so, since, without having definitely sold  it, I have touched the money.>

<Pay yourself,> added the unknown.

<I will do so, monsieur, since you so positively require it.>

A sad smile passed over the lips of the gentleman.

<Place the money on that trunk,> said he, turning round and  pointing to the piece of furniture.

Cropole deposited a tolerably large bag as directed, after having taken from it  the amount of his reckoning.

<Now,> said he, <I hope monsieur will not give me the pain of  not taking any supper. Dinner has already been refused; this is affronting to  the house of les Medici. Look, monsieur, the supper is on the table, and I  venture to say that it is not a bad one.>

The unknown asked for a glass of wine, broke off a morsel of bread, and did not  stir from the window whilst he ate and drank.

Shortly after was heard a loud flourish of trumpets; cries arose in the  distance, a confused buzzing filled the lower part of the city, and the first  distinct sound that struck the ears of the stranger was the tramp of advancing  horses.

<The king! the king!> repeated a noisy and eager crowd.

<The king!> cried Cropole, abandoning his guest and his ideas of  delicacy, to satisfy his curiosity.

With Cropole were mingled, and jostled, on the staircase, Madame Cropole,  Pittrino, and the waiters and scullions.

The cortège advanced slowly, lighted by a thousand flambeaux, in the streets  and from the windows.

After a company of musketeers, a closely ranked troop of gentlemen, came the  litter of monsieur le cardinal, drawn like a carriage by four black horses. The  pages and people of the cardinal marched behind.

Next came the carriage of the queen-mother, with her maids of honour at the  doors, her gentlemen on horseback at both sides.

The king then appeared, mounted upon a splendid horse of Saxon breed, with a  flowing mane. The young prince exhibited, when bowing to some windows from  which issued the most animated acclamations, a noble and handsome countenance,  illuminated by the flambeaux of his pages.

By the side of the king, though a little in the rear, the Prince de Condé, M.  Dangeau, and twenty other courtiers, followed by their people and their  baggage, closed this veritably triumphant march. The pomp was of a military  character.

Some of the courtiers—the elder ones, for instance—wore travelling  dresses; but all the rest were clothed in warlike panoply. Many wore the gorget  and buff coat of the times of Henry IV and Louis XIII

When the king passed before him, the unknown, who had leant forward over the  balcony to obtain a better view, and who had concealed his face by leaning on  his arm, felt his heart swell and overflow with a bitter jealousy.

The noise of the trumpets excited him—the popular acclamations deafened  him: for a moment he allowed his reason to be absorbed in this flood of lights,  tumult, and brilliant images.

<He is a king!> murmured he, in an accent of despair.

Then, before he had recovered from his sombre reverie, all the noise, all the  splendour, had passed away. At the angle of the street there remained nothing  beneath the stranger but a few hoarse, discordant voices, shouting at intervals  <Vive le Roi!>

There remained likewise the six candles held by the inhabitants of the hostelry  des Medici; that is to say, two for Cropole, two for Pittrino, and one for each  scullion. Cropole never ceased repeating, <How good-looking the king is!  How strongly he resembles his illustrious father!>

<A handsome likeness!> said Pittrino.

<And what a lofty carriage he has!> added Madame Cropole, already  in promiscuous commentary with her neighbours of both sexes.

Cropole was feeding their gossip with his own personal remarks, without  observing that an old man on foot, but leading a small Irish horse by the  bridle, was endeavouring to penetrate the crowd of men and women which blocked  up the entrance to the Medici. But at that moment the voice of the stranger was  heard from the window.

<Make way, monsieur l'hotelier, to the entrance of your  house!>

Cropole turned around, and, on seeing the old man, cleared a passage for him.

The window was instantly closed.

Pittrino pointed out the way to the newly-arrived guest, who entered without  uttering a word.

The stranger waited for him on the landing; he opened his arms to the old man,  and led him to a seat.

<Oh, no, no, my lord!> said he. <Sit down in your  presence?—never!>

<Parry,> cried the gentleman, <I beg you will; you come from  England—you come so far. Ah! it is not for your age to undergo the  fatigues my service requires. Rest yourself.>

<I have my reply to give your lordship, in the first place.>

<Parry, I conjure you to tell me nothing; for if your news had been good,  you would not have begun in such a manner; you go about, which proves that the  news is bad.>

<My lord,> said the old man, <do not hasten to alarm  yourself; all is not lost, I hope. You must employ energy, but more  particularly resignation.>

<Parry,> said the young man, <I have reached this place  through a thousand snares and after a thousand difficulties; can you doubt my  energy? I have meditated this journey ten years, in spite of all counsels and  all obstacles—have you faith in my perseverance? I have this evening sold  the last of my father's diamonds; for I had nothing wherewith to pay for  my lodgings and my host was about to turn me out.>

Parry made a gesture of indignation, to which the young man replied by a  pressure of the hand and a smile.

<I have still two hundred and seventy-four pistoles left and I feel  myself rich. I do not despair, Parry; have you faith in my resignation?>

The old man raised his trembling hands towards heaven.

<Let me know,> said the stranger,—<disguise nothing  from me—what has happened?>

<My recital will be short, my lord; but in the name of Heaven do not  tremble so.>

<It is impatience, Parry. Come, what did the general say to you?>

<At first the general would not receive me.>

<He took you for a spy?>

<Yes, my lord; but I wrote him a letter.>

<Well?>

<He read it, and received me, my lord.>

<Did that letter thoroughly explain my position and my views?>

<Oh, yes!> said Parry, with a sad smile; <it painted your  very thoughts faithfully.>

<Well—then, Parry.>

<Then the general sent me back the letter by an aide-de-camp, informing  me that if I were found the next day within the circumscription of his command,  he would have me arrested.>

<Arrested!> murmured the young man. <What! arrest you, my  most faithful servant?>

<Yes, my lord.>

<And notwithstanding you had signed the name Parry?>

<To all my letters, my lord; and the aide-de-camp had known me at St  James's and at Whitehall, too,> added the old man with a sigh.

The young man leaned forward, thoughtful and sad.

<Ay, that's what he did before his people,> said he,  endeavouring to cheat himself with hopes. <But, privately—between  you and him—what did he do? Answer!>

<Alas! my lord, he sent to me four cavaliers, who gave me the horse with  which you just now saw me come back. These cavaliers conducted me, in great  haste, to the little port of Tenby, threw me, rather than embarked me, into a  little fishing-boat, about to sail for Brittany, and here I am.>

<Oh!> sighed the young man, clasping his neck convulsively with his  hand, and with a sob. <Parry, is that all?—is that all?>

<Yes, my lord; that is all.>

After this brief reply ensued a long interval of silence, broken only by the  convulsive beating of the heel of the young man on the floor.

The old man endeavoured to change the conversation; it was leading to thoughts  much too sinister.

<My lord,> said he, <what is the meaning of all the noise  which preceded me? What are these people crying <Vive le Roi!> for?  What king do they mean? and what are all these lights for?>

<Ah! Parry,> replied the young man ironically, <don't  you know that this is the King of France visiting his good city of Blois? All  these trumpets are his, all those gilded housings are his, all those gentlemen  wear swords that are his. His mother precedes him in a carriage magnificently  encrusted with silver and gold. Happy mother! His minister heaps up millions,  and conducts him to a rich bride. Then all these people rejoice; they love  their king, they hail him with their acclamations, and they cry, <Vive le  Roi! Vive le Roi!>>

<Well, well, my lord,> said Parry, more uneasy at the turn the  conversation had taken than at the other.

<You know,> resumed the unknown, <that my mother and my  sister, whilst all this is going on in honour of the King of France, have  neither money nor bread; you know that I myself shall be poor and degraded  within a fortnight, when all Europe will become acquainted with what you have  told me. Parry, are there not examples in which a man of my condition should  himself\longdash>

<My lord, in the name of Heaven\longdash>

<You are right, Parry; I am a coward, and if I do nothing for myself,  what will God do? No, no; I have two arms, Parry, and I have a sword.>  And he struck his arm violently with his hand, and took down his sword, which  hung against the wall.

<What are you going to do, my lord?>

<What am I going to do, Parry? What every one in my family does. My  mother lives on public charity, my sister begs for my mother; I have, somewhere  or other, brothers who equally beg for themselves; and I, the eldest, will go  and do as all the rest do—I will go and ask charity!>

And with these words, which he finished sharply with a nervous and terrible  laugh, the young man girded on his sword, took his hat from the trunk, fastened  to his shoulder a black cloak, which he had worn all during his journey, and  pressing the two hands of the old man, who watched his proceedings with a look  of anxiety,—

<My good Parry,> said he, <order a fire, drink, eat, sleep,  and be happy; let us both be happy, my faithful friend, my only friend. We are  rich, as rich as kings!>

He struck the bag of pistoles with his clenched hand as he spoke, and it fell  heavily to the ground. He resumed that dismal laugh that had so alarmed Parry;  and whilst the whole household was screaming, singing, and preparing to install  the travellers who had been preceded by their lackeys, he glided out by the  principal entrance into the street, where the old man, who had gone to the  window, lost sight of him in a moment.  